\documentclass{article}
\usepackage{amsmath}
\usepackage{amssymb}

\title{A Compendium of Elementary Multivariable Calculus and Linear Algebra}
\author{Hyunseung Ryu}

\begin{document}
\maketitle

\begin{abstract}
This is a rendering test entry, not a real paper. It strings together standard
identities from multivariable calculus and linear algebra
\cite{apostol1969,strang2016} purely to exercise equation, matrix, and
citation rendering.
\end{abstract}

\section{Differential calculus in several variables}

For $f : \mathbb{R}^n \to \mathbb{R}$, the gradient and directional derivative
along a unit vector $\mathbf{u}$ are

\begin{align}
  \nabla f(\mathbf{x}) &= \left( \frac{\partial f}{\partial x_1}, \ldots,
    \frac{\partial f}{\partial x_n} \right), \\
  D_{\mathbf{u}} f(\mathbf{x}) &= \nabla f(\mathbf{x}) \cdot \mathbf{u}.
\end{align}

For $\mathbf{f} : \mathbb{R}^n \to \mathbb{R}^m$, the Jacobian is the matrix of
first partials,

\begin{equation}
  J_{\mathbf{f}}(\mathbf{x}) =
  \begin{bmatrix}
    \dfrac{\partial f_1}{\partial x_1} & \cdots & \dfrac{\partial f_1}{\partial x_n} \\
    \vdots & \ddots & \vdots \\
    \dfrac{\partial f_m}{\partial x_1} & \cdots & \dfrac{\partial f_m}{\partial x_n}
  \end{bmatrix},
\end{equation}

and the multivariable chain rule for $g(t) = f(\mathbf{x}(t))$ reads

\begin{equation}
  \frac{dg}{dt} = \sum_{i=1}^{n} \frac{\partial f}{\partial x_i} \frac{dx_i}{dt}
    = \nabla f(\mathbf{x}(t)) \cdot \mathbf{x}'(t).
\end{equation}

The Hessian collects second partials,

\begin{equation}
  H_f(\mathbf{x}) =
  \begin{bmatrix}
    \dfrac{\partial^2 f}{\partial x_1^2} & \cdots & \dfrac{\partial^2 f}{\partial x_1 \partial x_n} \\
    \vdots & \ddots & \vdots \\
    \dfrac{\partial^2 f}{\partial x_n \partial x_1} & \cdots & \dfrac{\partial^2 f}{\partial x_n^2}
  \end{bmatrix},
\end{equation}

and the second-order Taylor expansion about $\mathbf{a}$ is

\begin{equation}
  f(\mathbf{x}) \approx f(\mathbf{a}) + \nabla f(\mathbf{a})^{\!\top} (\mathbf{x} - \mathbf{a})
    + \frac{1}{2} (\mathbf{x} - \mathbf{a})^{\!\top} H_f(\mathbf{a}) (\mathbf{x} - \mathbf{a}).
\end{equation}

\section{Vector calculus}

For a vector field $\mathbf{F} = (P, Q, R)$, divergence and curl are

\begin{align}
  \nabla \cdot \mathbf{F} &= \frac{\partial P}{\partial x} + \frac{\partial Q}{\partial y}
    + \frac{\partial R}{\partial z}, \\
  \nabla \times \mathbf{F} &=
  \begin{vmatrix}
    \mathbf{i} & \mathbf{j} & \mathbf{k} \\
    \partial_x & \partial_y & \partial_z \\
    P & Q & R
  \end{vmatrix}.
\end{align}

The three classical integral theorems relate a boundary integral to an
integral over the enclosed region:

\begin{align}
  \oint_{\partial D} (P \, dx + Q \, dy) &= \iint_D \left( \frac{\partial Q}{\partial x}
    - \frac{\partial P}{\partial y} \right) dA
    && \text{(Green)}, \\
  \iint_S (\nabla \times \mathbf{F}) \cdot d\mathbf{S} &= \oint_{\partial S} \mathbf{F} \cdot d\mathbf{r}
    && \text{(Stokes)}, \\
  \iiint_V (\nabla \cdot \mathbf{F}) \, dV &= \oiint_{\partial V} \mathbf{F} \cdot d\mathbf{S}
    && \text{(Divergence)}.
\end{align}

\section{Linear algebra}

For $A \in \mathbb{R}^{n \times n}$, eigenpairs $(\lambda, \mathbf{v})$ satisfy

\begin{equation}
  A \mathbf{v} = \lambda \mathbf{v}, \qquad \det(A - \lambda I) = 0.
\end{equation}

If $A$ is diagonalizable, $A = P D P^{-1}$ with $D = \operatorname{diag}(\lambda_1,
\ldots, \lambda_n)$, and every real matrix admits a singular value decomposition

\begin{equation}
  A = U \Sigma V^{\!\top}, \qquad U^{\!\top} U = I, \quad V^{\!\top} V = I,
  \qquad \Sigma = \operatorname{diag}(\sigma_1, \ldots, \sigma_r, 0, \ldots).
\end{equation}

A symmetric matrix $A$ is positive definite iff its quadratic form is strictly
positive off the origin,

\begin{equation}
  \mathbf{x}^{\!\top} A \mathbf{x} > 0 \quad \forall \, \mathbf{x} \neq \mathbf{0}
  \iff \lambda_i > 0 \ \text{for all } i.
\end{equation}

For an inner product space, the Cauchy--Schwarz inequality bounds the inner
product by the norms,

\begin{equation}
  |\langle \mathbf{x}, \mathbf{y} \rangle| \le \lVert \mathbf{x} \rVert \, \lVert \mathbf{y} \rVert,
  \qquad \lVert \mathbf{x} \rVert = \sqrt{\langle \mathbf{x}, \mathbf{x} \rangle},
\end{equation}

and the orthogonal projection of $\mathbf{y}$ onto $\mathbf{x}$ is

\begin{equation}
  \operatorname{proj}_{\mathbf{x}} \mathbf{y}
    = \frac{\langle \mathbf{x}, \mathbf{y} \rangle}{\langle \mathbf{x}, \mathbf{x} \rangle} \, \mathbf{x}.
\end{equation}

\section{Constrained optimization}

For $f, g : \mathbb{R}^n \to \mathbb{R}$, a critical point of $f$ subject to
$g(\mathbf{x}) = 0$ satisfies, for some multiplier $\lambda$,

\begin{equation}
  \nabla f(\mathbf{x}^{*}) = \lambda \, \nabla g(\mathbf{x}^{*}), \qquad
  g(\mathbf{x}^{*}) = 0,
\end{equation}

equivalently $\nabla \mathcal{L} = \mathbf{0}$ for the Lagrangian

\begin{equation}
  \mathcal{L}(\mathbf{x}, \lambda) = f(\mathbf{x}) - \lambda \, g(\mathbf{x}).
\end{equation}

\section{Conclusion}

Every identity above is standard \cite{strang2016,apostol1969}; the entry
exists to check that display equations, aligned systems, matrices, and
citations all render correctly, not to report a result.

\end{document}
